\tzpanchor{1268}
\tzpentryheader{1268}{VON NEUMANN-LINDBLAD DYNAMICS}

\begin{tzpsnapshot}
\tzpsnaprow{TZPID}{\tzpcode{ID1268}}
\tzpsnaprow{UUID}{\tzpcode{538d861c-3df7-5b1e-b9e5-c5f0189f4e63}}
\tzpsnaprow{Type}{Transfer Statement}
\tzpsnaprow{Scale}{transfer}
\tzpsnaprow{LOC Classification}{Working routing: QA641 manifolds and topology; QC174.17 mathematical physics}
\tzpsnaprow{arXiv Classification}{Primary: math-ph; Cross-list: gr-qc, math.DG}
\end{tzpsnapshot}

\tzpsection{Canonical Statement}
2. Phi(g o f) = Phi(g) o Phi(f) (composition preservation)

\[
Phi(g ∘ f) = \Phi(g) ∘ \Phi(f) (composition preservation)
\]
\[
mathcal{J}_{\mathrm{out}} = \mathcal{T}_{\mathrm{TZP}}\!\left(\mathcal{K}_{\mathrm{TZP}}\right)
\]

\noindent
\begin{minipage}[t]{0.48\linewidth}
\tzpsection{Wolfram Form}
\begingroup
\small
\[
\begin{adjustbox}{max width=\linewidth}
$\displaystyle
\mathfrak{W}_{\mathrm{TZP}}\!\left(\mathcal{K}_{1268}\right)
\longmapsto
\operatorname{WeakProtocolKey}\!\left\langle \mathcal{E}_{\mathrm{decoherence}},\; K_{k},\; p_{\mathrm{scatter}},\; p_{\mathrm{abs}} \right\rangle
$
\end{adjustbox}
\]
\[
\begin{adjustbox}{max width=\linewidth}
$\displaystyle
\operatorname{ProtocolSource}\!\left(\mathcal{K}_{1268}\right)
=
\operatorname{HoldForm}\!\left(\mathrm{decoherence superoperator}\right)
$
\end{adjustbox}
\]
\textbf{Wolfram register:} weak-value protocol formalism lift\\
\textbf{Title lift:} VON NEUMANN-LINDBLAD DYNAMICS\\
\textbf{Protocol source:} Decoherence superoperator\\
\textbf{Symbol key:} \(\mathcal{E}_{\mathrm{decoherence}}\): decoherence channel acting on the input state\\ \(K_{k}\): Kraus operator for the environmental channel\\ \(p_{\mathrm{scatter}}\): scattering probability in the noisy channel\\ \(p_{\mathrm{abs}}\): absorption probability into the vacuum state
\par
\endgroup
\end{minipage}\hfill
\begin{minipage}[t]{0.48\linewidth}
\tzpsection{TRAWIN Normalization}
\[
\tzpnormalizewrap{Phi(g ∘ f) = \Phi(g) ∘ \Phi(f) (composition preservation),\,mathcal{J}_{\mathrm{out}} = \mathcal{T}_{\mathrm{TZP}}\!\left(\mathcal{K}_{\mathrm{TZP}}\right),\,mathcal{A}_{\mathrm{adm}}\!\left(\mathcal{J}_{\mathrm{out}}\right)=1}
\]

\small
\textbf{T}: Canonical Definition is localized within the engineering and implementation arcs arc of the directory.\\
\textbf{R}: Support Relation preserves the operative relation that von neumann-lindblad dynamics requires.\\
\textbf{A}: Closure Relation fixes the admissible closure condition for the entry.\\
\textbf{W}: VON NEUMANN-LINDBLAD DYNAMICS is rendered as a reusable operator-level statement.\\
\textbf{I}: The informational content of von neumann-lindblad dynamics is retained rather than flattened.\\
\textbf{N}: VON NEUMANN-LINDBLAD DYNAMICS closes as a distinct normalized TRAWIN object.
\normalsize
\end{minipage}

\tzpsection{Derivable Consequences}
The entry now reads through actual sourced equations, so its local arc is carried by explicit mathematical relations rather than a uniform quaternary certificate record shell.
\clearpage

\tzpentryheader{1268}{Oxford Dictionary of the Queens, NY English}

\begingroup\small
\tzpsection{Definition}
VON NEUMANN-LINDBLAD DYNAMICS is a registry definition in Engineering and Implementation Arcs with secondary pressure from Registry and Publication Workflow.

% BEGIN TZP CONTEXTUAL MEANING
\tzpsection{Contextual Meaning}
\begingroup\footnotesize\setlength{\emergencystretch}{3em}\sloppy
\textbf{Formal statement:} 2. Phi(g o f) = Phi(g) o Phi(f) (composition preservation). \textbf{Contextual meaning:} The entry reads VON NEUMANN-LINDBLAD DYNAMICS as a controlled technical headword whose equation layer, proof route, and registry role are interpreted together. \textbf{Derivable consequences:} The entry now reads through actual sourced equations, so its local arc is carried by explicit mathematical relations rather than a uniform quaternary certificate record shell. \textbf{Lemma support:} Canonical Definition; Support Relation; Closure Relation.\par
\endgroup
% END TZP CONTEXTUAL MEANING

\tzpsection{Technical Interpretation}
Technically, VON NEUMANN-LINDBLAD DYNAMICS is read through the local equation chain and the TRAWIN normalization recorded on the entry page.

\tzpsection{Equation Grounding}
Equation anchors: Phi(g   f) = Phi(g)   Phi(f) (composition preservation); mathcal J sub out = T sub TZP ( K sub TZP ); and mathcal A sub adm ( J sub out )=1. Proof route: show that the stated equations form a coherent progression for the entry and remain compatible under the TRAWIN ordering. Formalization route: prove that the final closure relation follows from the preceding entry-specific equations.

\tzpsection{Interpretive Role}
The entry now reads through actual sourced equations, so its local arc is carried by explicit mathematical relations rather than a uniform quaternary certificate record shell.
\endgroup

\tzpsection{TZP Thesaurus}
\begin{enumerate}[leftmargin=1.6em,itemsep=6pt,topsep=4pt]
\item \textbf{Von Neumann Lindblad Dynamics}: entry-specific alternate wording for indexing, related literature, and local formal context.
\item \textbf{Von Neumann Lindblad Formal Analogue}: entry-specific alternate wording for indexing, related literature, and local formal context.
\item \textbf{Dynamics Equivalence}: entry-specific alternate wording for indexing, related literature, and local formal context.
\end{enumerate}

\tzpsection{Encyclopedia Mathematica}
\begin{samepage}
\noindent\begin{minipage}[t]{0.45\linewidth}
\vspace{0pt}
\raggedright
\scriptsize\textbf{Image reading.} The rendered manifold at right is the per-ID light plate for VON NEUMANN-LINDBLAD DYNAMICS. It should be read as the visual companion to the entry's equation layer, not as a repeated template diagram.\par
\end{minipage}\hfill
\begin{minipage}[t]{0.53\linewidth}
\vspace*{-0.10in}
\raggedright
\tzpencyclopediaimage{1268}
\end{minipage}
\end{samepage}

\clearpage

\tzpentryheader{1268}{Direct Equation Progression and Proof Trajectory}

\tzpsection{Direct Equation Progression}
\begin{enumerate}[leftmargin=1.6em,itemsep=9pt,topsep=6pt]
\item \textbf{Canonical Definition}
\[
Phi(g ∘ f) = \Phi(g) ∘ \Phi(f) (composition preservation)
\]
This fixes the first explicit relation carried by the entry rather than a generic certificate record normalization.

\item \textbf{Support Relation}
\[
mathcal{J}_{\mathrm{out}} = \mathcal{T}_{\mathrm{TZP}}\!\left(\mathcal{K}_{\mathrm{TZP}}\right)
\]
This adds the next operational equation that advances the entry beyond its opening definition.

\item \textbf{Closure Relation}
\[
mathcal{A}_{\mathrm{adm}}\!\left(\mathcal{J}_{\mathrm{out}}\right)=1
\]
This closes the progression with an actual admissibility or closure relation instead of recycling the normalization shell.
\end{enumerate}

\tzpsection{Proof}
\textbf{Proof route.} show that the stated equations form a coherent progression for the entry and remain compatible under the TRAWIN ordering.

\textbf{Formal method.} prove that the final closure relation follows from the preceding entry-specific equations.

\medskip
\tzpproofartifacts{Isabelle audit: corpus classification recorded in appendix.}{Lean audit: compiled corpus certificate.}{Rocq audit: compiled corpus certificate.}

\tzpsection{Dependencies and Test Handle}
\begin{tzpsnapshot}
\tzpsnaprow{Dependencies}{\tzpidref{1267}, \tzpidref{1200}}
\tzpsnaprow{Node}{\tzpcode{registry_id_1268}}
\tzpsnaprow{Support Titles}{Canonical Definition; Support Relation; Closure Relation}
\tzpsnaprow{Test Handle}{If the cited entry equations cannot be made mutually admissible in one TRAWIN-normalized frame, the entry fails.}
\end{tzpsnapshot}

\tzpsection{Canonical Equation Links}
\tzpequationlinks
  {Phi(g ∘ f) = \Phi(g) ∘ \Phi(f) (composition preservation)}
  {\mathcal{J}_{\mathrm{out}} = \cdots}
  {\mathcal{A}_{\mathrm{adm}}\!\left(\mathcal{J}_{\mathrm{out}}\right) = \cdots}
